% --- Archon physics formalization source begin ---
% archon:physics
% archon:covers PhyXMiniProblems/problem_phyx_mini_0664.lean
% archon:source-report reports/phyx_mini/problem_phyx_mini_0664.source.json

\chapter{Physics problem phyx\_mini\_0664}
\label{ch:PhyXMiniProblems_problem_phyx_mini_0664}

\paragraph{Problem source.}
\#\# Physical scenario

A steel tank of cross-sectional area \$3 \textbackslash{}, \textbackslash{}text\{m\}\textasciicircum{}2\$ and height \$16 \textbackslash{}, \textbackslash{}text\{m\}\$ weighs \$10\textbackslash{},000 \textbackslash{}, \textbackslash{}text\{kg\}\$ and is open at the top, as shown in figure. We want to float it in the ocean so that it is positioned \$10 \textbackslash{}, \textbackslash{}text\{m\}\$ straight down by pouring concrete into its bottom.

\#\# Question

How much concrete should we use?

\#\# Answer choices

- A: A: 10000 kg

- B: B: 29910 kg

- C: C: 9970 kg

- D: D: 19910 kg

\#\# Figure caption (auxiliary; use the image as primary evidence)

The image is a diagram depicting a cross-sectional view of a structure submerged in the ocean. The structure appears to be a vertical cylinder or tank with an open top above the water surface.

- The structure has three labeled sections: "Air," "Concrete," and "Ocean."
- The top section is labeled "Air," indicating the upper part of the structure is filled with air.
- The bottom section is labeled "Concrete," indicating that this part is filled with concrete.
- The "Ocean" is labeled outside the structure, representing the surrounding water.
- The water level in the ocean is at the same height as the top of the air section inside the structure.
- There is a double-headed vertical arrow on the left side labeled "10 m," indicating that the height from the base of the concrete section to the ocean surface is 10 meters.
- The ocean's surface is represented by a wavy line. 

The diagram illustrates the relationship between the structure's internal components and its submersion in the ocean.

\#\# Dataset metadata

Statics; Physical Model Grounding Reasoning, Multi-Formula Reasoning

\paragraph{Recorded answer/context.}
D: D: 19910 kg

\paragraph{Figure/image path.}
/root/proposal\_for\_physic/science-mango/phyx\_mini\_run/phyx\_data/test\_image/664.png

\paragraph{Formalization target.}
create a compiling Lean file with sorry bodies at `PhyXMiniProblems/problem\_phyx\_mini\_0664.lean`. The Lean declarations must preserve the physical quantities, dimensions or dimensional roles, figure labels, governing-law hypotheses, and final relation expressed by this problem.
Use Mathlib/Physlib names found through LeanExplore where available. If a physics API is missing, introduce faithful local abstractions rather than scalar placeholder aliases.

\begin{theorem}[Physics formalization target]
\label{thm:physics:phyx_mini_0664:target}
\lean{PhyXMiniProblems.ProblemPhyXMini0664.problem_phyx_mini_0664}
\uses{def:physics:phyx-mini-0664:phyxminiproblems-problemphyxmini0664-massinkilograms, def:physics:phyx-mini-0664:phyxminiproblems-problemphyxmini0664-ballastedtanksetup, def:physics:phyx-mini-0664:phyxminiproblems-problemphyxmini0664-matchesproblemandfiguredata, def:physics:phyx-mini-0664:phyxminiproblems-problemphyxmini0664-usesrecordedoceandensitycalibration, def:physics:phyx-mini-0664:phyxminiproblems-problemphyxmini0664-hasphysicalgravity, def:physics:phyx-mini-0664:phyxminiproblems-problemphyxmini0664-satisfiesprismaticdisplacementlaw, def:physics:phyx-mini-0664:phyxminiproblems-problemphyxmini0664-satisfiestotalweightlaw, def:physics:phyx-mini-0664:phyxminiproblems-problemphyxmini0664-satisfiesarchimedesprinciple, def:physics:phyx-mini-0664:phyxminiproblems-problemphyxmini0664-floatsinverticalstaticequilibrium, def:physics:phyx-mini-0664:phyxminiproblems-problemphyxmini0664-matchesanswerchoice}
The assigned autoformalize agent should translate this physics problem into Lean declarations in the covered file, with theorem and lemma proof bodies written as `by sorry`.
\end{theorem}
\begin{proof}
This is an autoformalization task, not a proof task. Produce statements that can later be proved by the physics prover without weakening the physical model.
\end{proof}

% --- Archon declaration topology begin ---
\section{Declaration topology}
\begin{definition}[Length Quantity]
  \label{def:physics:phyx-mini-0664:phyxminiproblems-problemphyxmini0664-lengthquantity}
  \lean{PhyXMiniProblems.ProblemPhyXMini0664.LengthQuantity}
  A physical length
\end{definition}
\begin{proof}
  This is the declared model definition of Length Quantity.
\end{proof}
\begin{definition}[Area Quantity]
  \label{def:physics:phyx-mini-0664:phyxminiproblems-problemphyxmini0664-areaquantity}
  \lean{PhyXMiniProblems.ProblemPhyXMini0664.AreaQuantity}
  A physical cross-sectional area
\end{definition}
\begin{proof}
  This is the declared model definition of Area Quantity.
\end{proof}
\begin{definition}[Volume Quantity]
  \label{def:physics:phyx-mini-0664:phyxminiproblems-problemphyxmini0664-volumequantity}
  \lean{PhyXMiniProblems.ProblemPhyXMini0664.VolumeQuantity}
  A physical volume
\end{definition}
\begin{proof}
  This is the declared model definition of Volume Quantity.
\end{proof}
\begin{definition}[Mass Quantity]
  \label{def:physics:phyx-mini-0664:phyxminiproblems-problemphyxmini0664-massquantity}
  \lean{PhyXMiniProblems.ProblemPhyXMini0664.MassQuantity}
  A physical mass
\end{definition}
\begin{proof}
  This is the declared model definition of Mass Quantity.
\end{proof}
\begin{definition}[Mass Density Quantity]
  \label{def:physics:phyx-mini-0664:phyxminiproblems-problemphyxmini0664-massdensityquantity}
  \lean{PhyXMiniProblems.ProblemPhyXMini0664.MassDensityQuantity}
  A physical mass density, with dimension M L⁻³
\end{definition}
\begin{proof}
  This is the declared model definition of Mass Density Quantity.
\end{proof}
\begin{definition}[Acceleration Quantity]
  \label{def:physics:phyx-mini-0664:phyxminiproblems-problemphyxmini0664-accelerationquantity}
  \lean{PhyXMiniProblems.ProblemPhyXMini0664.AccelerationQuantity}
  A physical acceleration magnitude
\end{definition}
\begin{proof}
  This is the declared model definition of Acceleration Quantity.
\end{proof}
\begin{definition}[Force Quantity]
  \label{def:physics:phyx-mini-0664:phyxminiproblems-problemphyxmini0664-forcequantity}
  \lean{PhyXMiniProblems.ProblemPhyXMini0664.ForceQuantity}
  A physical force magnitude
\end{definition}
\begin{proof}
  This is the declared model definition of Force Quantity.
\end{proof}
\begin{definition}[si Readout]
  \label{def:physics:phyx-mini-0664:phyxminiproblems-problemphyxmini0664-sireadout}
  \lean{PhyXMiniProblems.ProblemPhyXMini0664.siReadout}
  Read a dimensionful quantity in the coherent SI unit system
\end{definition}
\begin{proof}
  This is the declared model definition of si Readout.
\end{proof}
\begin{definition}[length In Meters]
  \label{def:physics:phyx-mini-0664:phyxminiproblems-problemphyxmini0664-lengthinmeters}
  \lean{PhyXMiniProblems.ProblemPhyXMini0664.lengthInMeters}
  \uses{def:physics:phyx-mini-0664:phyxminiproblems-problemphyxmini0664-lengthquantity, def:physics:phyx-mini-0664:phyxminiproblems-problemphyxmini0664-sireadout}
  Metre readout of a physical length
\end{definition}
\begin{proof}
  This is the declared model definition of length In Meters.
\end{proof}
\begin{definition}[area In Square Meters]
  \label{def:physics:phyx-mini-0664:phyxminiproblems-problemphyxmini0664-areainsquaremeters}
  \lean{PhyXMiniProblems.ProblemPhyXMini0664.areaInSquareMeters}
  \uses{def:physics:phyx-mini-0664:phyxminiproblems-problemphyxmini0664-areaquantity, def:physics:phyx-mini-0664:phyxminiproblems-problemphyxmini0664-sireadout}
  Square-metre readout of a physical area
\end{definition}
\begin{proof}
  This is the declared model definition of area In Square Meters.
\end{proof}
\begin{definition}[volume In Cubic Meters]
  \label{def:physics:phyx-mini-0664:phyxminiproblems-problemphyxmini0664-volumeincubicmeters}
  \lean{PhyXMiniProblems.ProblemPhyXMini0664.volumeInCubicMeters}
  \uses{def:physics:phyx-mini-0664:phyxminiproblems-problemphyxmini0664-volumequantity, def:physics:phyx-mini-0664:phyxminiproblems-problemphyxmini0664-sireadout}
  Cubic-metre readout of a physical volume
\end{definition}
\begin{proof}
  This is the declared model definition of volume In Cubic Meters.
\end{proof}
\begin{definition}[mass In Kilograms]
  \label{def:physics:phyx-mini-0664:phyxminiproblems-problemphyxmini0664-massinkilograms}
  \lean{PhyXMiniProblems.ProblemPhyXMini0664.massInKilograms}
  \uses{def:physics:phyx-mini-0664:phyxminiproblems-problemphyxmini0664-massquantity, def:physics:phyx-mini-0664:phyxminiproblems-problemphyxmini0664-sireadout}
  Kilogram readout of a physical mass
\end{definition}
\begin{proof}
  This is the declared model definition of mass In Kilograms.
\end{proof}
\begin{definition}[density In Kilograms Per Cubic Meter]
  \label{def:physics:phyx-mini-0664:phyxminiproblems-problemphyxmini0664-densityinkilogramspercubicmeter}
  \lean{PhyXMiniProblems.ProblemPhyXMini0664.densityInKilogramsPerCubicMeter}
  \uses{def:physics:phyx-mini-0664:phyxminiproblems-problemphyxmini0664-massdensityquantity, def:physics:phyx-mini-0664:phyxminiproblems-problemphyxmini0664-sireadout}
  Kilogram-per-cubic-metre readout of a mass density
\end{definition}
\begin{proof}
  This is the declared model definition of density In Kilograms Per Cubic Meter.
\end{proof}
\begin{definition}[acceleration In Meters Per Second Squared]
  \label{def:physics:phyx-mini-0664:phyxminiproblems-problemphyxmini0664-accelerationinmeterspersecondsquared}
  \lean{PhyXMiniProblems.ProblemPhyXMini0664.accelerationInMetersPerSecondSquared}
  \uses{def:physics:phyx-mini-0664:phyxminiproblems-problemphyxmini0664-accelerationquantity, def:physics:phyx-mini-0664:phyxminiproblems-problemphyxmini0664-sireadout}
  Metre-per-second-squared readout of an acceleration magnitude
\end{definition}
\begin{proof}
  This is the declared model definition of acceleration In Meters Per Second Squared.
\end{proof}
\begin{definition}[force In Newtons]
  \label{def:physics:phyx-mini-0664:phyxminiproblems-problemphyxmini0664-forceinnewtons}
  \lean{PhyXMiniProblems.ProblemPhyXMini0664.forceInNewtons}
  \uses{def:physics:phyx-mini-0664:phyxminiproblems-problemphyxmini0664-forcequantity, def:physics:phyx-mini-0664:phyxminiproblems-problemphyxmini0664-sireadout}
  Newton readout of a force magnitude
\end{definition}
\begin{proof}
  This is the declared model definition of force In Newtons.
\end{proof}
\begin{definition}[Tank Top Boundary]
  \label{def:physics:phyx-mini-0664:phyxminiproblems-problemphyxmini0664-tanktopboundary}
  \lean{PhyXMiniProblems.ProblemPhyXMini0664.TankTopBoundary}
  Whether the tank's top boundary is open or sealed
\end{definition}
\begin{proof}
  This is the declared model definition of Tank Top Boundary.
\end{proof}
\begin{definition}[Tank Axis Orientation]
  \label{def:physics:phyx-mini-0664:phyxminiproblems-problemphyxmini0664-tankaxisorientation}
  \lean{PhyXMiniProblems.ProblemPhyXMini0664.TankAxisOrientation}
  Orientation of the long axis of the tank
\end{definition}
\begin{proof}
  This is the declared model definition of Tank Axis Orientation.
\end{proof}
\begin{definition}[Tank Shell Material]
  \label{def:physics:phyx-mini-0664:phyxminiproblems-problemphyxmini0664-tankshellmaterial}
  \lean{PhyXMiniProblems.ProblemPhyXMini0664.TankShellMaterial}
  Material of the tank shell
\end{definition}
\begin{proof}
  This is the declared model definition of Tank Shell Material.
\end{proof}
\begin{definition}[Ballast Location]
  \label{def:physics:phyx-mini-0664:phyxminiproblems-problemphyxmini0664-ballastlocation}
  \lean{PhyXMiniProblems.ProblemPhyXMini0664.BallastLocation}
  Location of the concrete ballast inside the tank
\end{definition}
\begin{proof}
  This is the declared model definition of Ballast Location.
\end{proof}
\begin{definition}[Figure Region Label]
  \label{def:physics:phyx-mini-0664:phyxminiproblems-problemphyxmini0664-figureregionlabel}
  \lean{PhyXMiniProblems.ProblemPhyXMini0664.FigureRegionLabel}
  Text labels visible in the supplied cross-sectional figure
\end{definition}
\begin{proof}
  This is the declared model definition of Figure Region Label.
\end{proof}
\begin{definition}[Supplied Figure Readout]
  \label{def:physics:phyx-mini-0664:phyxminiproblems-problemphyxmini0664-suppliedfigurereadout}
  \lean{PhyXMiniProblems.ProblemPhyXMini0664.SuppliedFigureReadout}
  \uses{def:physics:phyx-mini-0664:phyxminiproblems-problemphyxmini0664-lengthquantity, def:physics:phyx-mini-0664:phyxminiproblems-problemphyxmini0664-figureregionlabel}
  Primary-image readout. No numerical concrete height is introduced because the figure does not label one
\end{definition}
\begin{proof}
  This is the declared model definition of Supplied Figure Readout.
\end{proof}
\begin{definition}[Ballasted Tank Setup]
  \label{def:physics:phyx-mini-0664:phyxminiproblems-problemphyxmini0664-ballastedtanksetup}
  \lean{PhyXMiniProblems.ProblemPhyXMini0664.BallastedTankSetup}
  \uses{def:physics:phyx-mini-0664:phyxminiproblems-problemphyxmini0664-lengthquantity, def:physics:phyx-mini-0664:phyxminiproblems-problemphyxmini0664-areaquantity, def:physics:phyx-mini-0664:phyxminiproblems-problemphyxmini0664-volumequantity, def:physics:phyx-mini-0664:phyxminiproblems-problemphyxmini0664-massquantity, def:physics:phyx-mini-0664:phyxminiproblems-problemphyxmini0664-massdensityquantity, def:physics:phyx-mini-0664:phyxminiproblems-problemphyxmini0664-accelerationquantity, def:physics:phyx-mini-0664:phyxminiproblems-problemphyxmini0664-forcequantity, def:physics:phyx-mini-0664:phyxminiproblems-problemphyxmini0664-tanktopboundary, def:physics:phyx-mini-0664:phyxminiproblems-problemphyxmini0664-tankaxisorientation, def:physics:phyx-mini-0664:phyxminiproblems-problemphyxmini0664-tankshellmaterial, def:physics:phyx-mini-0664:phyxminiproblems-problemphyxmini0664-ballastlocation, def:physics:phyx-mini-0664:phyxminiproblems-problemphyxmini0664-suppliedfigurereadout}
  Physical state of the floating tank. In particular, concreteMass is an independent unknown and is not defined using the requested answer
\end{definition}
\begin{proof}
  This is the declared model definition of Ballasted Tank Setup.
\end{proof}
\begin{definition}[draft In Meters]
  \label{def:physics:phyx-mini-0664:phyxminiproblems-problemphyxmini0664-draftinmeters}
  \lean{PhyXMiniProblems.ProblemPhyXMini0664.draftInMeters}
  \uses{def:physics:phyx-mini-0664:phyxminiproblems-problemphyxmini0664-lengthinmeters, def:physics:phyx-mini-0664:phyxminiproblems-problemphyxmini0664-ballastedtanksetup}
  The draft indicated by the figure's vertical 10 m arrow
\end{definition}
\begin{proof}
  This is the declared model definition of draft In Meters.
\end{proof}
\begin{definition}[total Supported Mass In Kilograms]
  \label{def:physics:phyx-mini-0664:phyxminiproblems-problemphyxmini0664-totalsupportedmassinkilograms}
  \lean{PhyXMiniProblems.ProblemPhyXMini0664.totalSupportedMassInKilograms}
  \uses{def:physics:phyx-mini-0664:phyxminiproblems-problemphyxmini0664-massinkilograms, def:physics:phyx-mini-0664:phyxminiproblems-problemphyxmini0664-ballastedtanksetup}
  Total tank-plus-ballast mass readout, neglecting the contained air mass
\end{definition}
\begin{proof}
  This is the declared model definition of total Supported Mass In Kilograms.
\end{proof}
\begin{definition}[Matches Problem And Figure Data]
  \label{def:physics:phyx-mini-0664:phyxminiproblems-problemphyxmini0664-matchesproblemandfiguredata}
  \lean{PhyXMiniProblems.ProblemPhyXMini0664.MatchesProblemAndFigureData}
  \uses{def:physics:phyx-mini-0664:phyxminiproblems-problemphyxmini0664-lengthinmeters, def:physics:phyx-mini-0664:phyxminiproblems-problemphyxmini0664-areainsquaremeters, def:physics:phyx-mini-0664:phyxminiproblems-problemphyxmini0664-massinkilograms, def:physics:phyx-mini-0664:phyxminiproblems-problemphyxmini0664-ballastedtanksetup, def:physics:phyx-mini-0664:phyxminiproblems-problemphyxmini0664-draftinmeters}
  Data stated in the prose or read from the primary figure. The tank is dry inside, has air over concrete, and extends 16 m from bottom to its open top; the ocean surface is 10 m above its bottom
\end{definition}
\begin{proof}
  This is the declared model definition of Matches Problem And Figure Data.
\end{proof}
\begin{definition}[Uses Recorded Ocean Density Calibration]
  \label{def:physics:phyx-mini-0664:phyxminiproblems-problemphyxmini0664-usesrecordedoceandensitycalibration}
  \lean{PhyXMiniProblems.ProblemPhyXMini0664.UsesRecordedOceanDensityCalibration}
  \uses{def:physics:phyx-mini-0664:phyxminiproblems-problemphyxmini0664-densityinkilogramspercubicmeter, def:physics:phyx-mini-0664:phyxminiproblems-problemphyxmini0664-ballastedtanksetup}
  Environmental calibration implicit in the recorded exact answer. The source calls the fluid ocean water but supplies no density; 997 kg/m³ is therefore kept as an explicit extra datum rather than hidden in Archimedes' law
\end{definition}
\begin{proof}
  This is the declared model definition of Uses Recorded Ocean Density Calibration.
\end{proof}
\begin{definition}[Has Physical Gravity]
  \label{def:physics:phyx-mini-0664:phyxminiproblems-problemphyxmini0664-hasphysicalgravity}
  \lean{PhyXMiniProblems.ProblemPhyXMini0664.HasPhysicalGravity}
  \uses{def:physics:phyx-mini-0664:phyxminiproblems-problemphyxmini0664-accelerationinmeterspersecondsquared, def:physics:phyx-mini-0664:phyxminiproblems-problemphyxmini0664-ballastedtanksetup}
  Positivity needed to cancel gravity in the later static-equilibrium proof
\end{definition}
\begin{proof}
  This is the declared model definition of Has Physical Gravity.
\end{proof}
\begin{definition}[Satisfies Prismatic Displacement Law]
  \label{def:physics:phyx-mini-0664:phyxminiproblems-problemphyxmini0664-satisfiesprismaticdisplacementlaw}
  \lean{PhyXMiniProblems.ProblemPhyXMini0664.SatisfiesPrismaticDisplacementLaw}
  \uses{def:physics:phyx-mini-0664:phyxminiproblems-problemphyxmini0664-areainsquaremeters, def:physics:phyx-mini-0664:phyxminiproblems-problemphyxmini0664-volumeincubicmeters, def:physics:phyx-mini-0664:phyxminiproblems-problemphyxmini0664-ballastedtanksetup, def:physics:phyx-mini-0664:phyxminiproblems-problemphyxmini0664-draftinmeters}
  Prismatic displacement law: for the upright unflooded tank, displaced ocean volume is the horizontal cross-sectional area times the draft
\end{definition}
\begin{proof}
  This is the declared model definition of Satisfies Prismatic Displacement Law.
\end{proof}
\begin{definition}[Satisfies Total Weight Law]
  \label{def:physics:phyx-mini-0664:phyxminiproblems-problemphyxmini0664-satisfiestotalweightlaw}
  \lean{PhyXMiniProblems.ProblemPhyXMini0664.SatisfiesTotalWeightLaw}
  \uses{def:physics:phyx-mini-0664:phyxminiproblems-problemphyxmini0664-accelerationinmeterspersecondsquared, def:physics:phyx-mini-0664:phyxminiproblems-problemphyxmini0664-forceinnewtons, def:physics:phyx-mini-0664:phyxminiproblems-problemphyxmini0664-ballastedtanksetup, def:physics:phyx-mini-0664:phyxminiproblems-problemphyxmini0664-totalsupportedmassinkilograms}
  The downward weight magnitude is (tank mass + concrete mass) g
\end{definition}
\begin{proof}
  This is the declared model definition of Satisfies Total Weight Law.
\end{proof}
\begin{definition}[Satisfies Archimedes Principle]
  \label{def:physics:phyx-mini-0664:phyxminiproblems-problemphyxmini0664-satisfiesarchimedesprinciple}
  \lean{PhyXMiniProblems.ProblemPhyXMini0664.SatisfiesArchimedesPrinciple}
  \uses{def:physics:phyx-mini-0664:phyxminiproblems-problemphyxmini0664-volumeincubicmeters, def:physics:phyx-mini-0664:phyxminiproblems-problemphyxmini0664-densityinkilogramspercubicmeter, def:physics:phyx-mini-0664:phyxminiproblems-problemphyxmini0664-accelerationinmeterspersecondsquared, def:physics:phyx-mini-0664:phyxminiproblems-problemphyxmini0664-forceinnewtons, def:physics:phyx-mini-0664:phyxminiproblems-problemphyxmini0664-ballastedtanksetup}
  Archimedes' law: buoyancy is ocean density × displaced volume × g
\end{definition}
\begin{proof}
  This is the declared model definition of Satisfies Archimedes Principle.
\end{proof}
\begin{definition}[Floats In Vertical Static Equilibrium]
  \label{def:physics:phyx-mini-0664:phyxminiproblems-problemphyxmini0664-floatsinverticalstaticequilibrium}
  \lean{PhyXMiniProblems.ProblemPhyXMini0664.FloatsInVerticalStaticEquilibrium}
  \uses{def:physics:phyx-mini-0664:phyxminiproblems-problemphyxmini0664-forceinnewtons, def:physics:phyx-mini-0664:phyxminiproblems-problemphyxmini0664-ballastedtanksetup}
  Static vertical flotation: upward buoyancy balances downward weight
\end{definition}
\begin{proof}
  This is the declared model definition of Floats In Vertical Static Equilibrium.
\end{proof}
\begin{definition}[Answer Choice]
  \label{def:physics:phyx-mini-0664:phyxminiproblems-problemphyxmini0664-answerchoice}
  \lean{PhyXMiniProblems.ProblemPhyXMini0664.AnswerChoice}
  The supplied answer-choice labels
\end{definition}
\begin{proof}
  This is the declared model definition of Answer Choice.
\end{proof}
\begin{definition}[mass In Kilograms]
  \label{def:physics:phyx-mini-0664:phyxminiproblems-problemphyxmini0664-answerchoice-massinkilograms}
  \lean{PhyXMiniProblems.ProblemPhyXMini0664.AnswerChoice.massInKilograms}
  \uses{def:physics:phyx-mini-0664:phyxminiproblems-problemphyxmini0664-massinkilograms, def:physics:phyx-mini-0664:phyxminiproblems-problemphyxmini0664-answerchoice}
  Concrete-mass readout printed beside each answer choice
\end{definition}
\begin{proof}
  This is the declared model definition of mass In Kilograms.
\end{proof}
\begin{definition}[Matches Answer Choice]
  \label{def:physics:phyx-mini-0664:phyxminiproblems-problemphyxmini0664-matchesanswerchoice}
  \lean{PhyXMiniProblems.ProblemPhyXMini0664.MatchesAnswerChoice}
  \uses{def:physics:phyx-mini-0664:phyxminiproblems-problemphyxmini0664-massinkilograms, def:physics:phyx-mini-0664:phyxminiproblems-problemphyxmini0664-ballastedtanksetup, def:physics:phyx-mini-0664:phyxminiproblems-problemphyxmini0664-answerchoice}
  A modeled concrete mass agrees with a displayed answer choice
\end{definition}
\begin{proof}
  This is the declared model definition of Matches Answer Choice.
\end{proof}
\begin{lemma}[displaced ocean volume at requested draft]
  \label{lem:physics:phyx-mini-0664:phyxminiproblems-problemphyxmini0664-displaced-ocean-volume-at-requested-draft}
  \lean{PhyXMiniProblems.ProblemPhyXMini0664.displaced_ocean_volume_at_requested_draft}
  \uses{def:physics:phyx-mini-0664:phyxminiproblems-problemphyxmini0664-volumeincubicmeters, def:physics:phyx-mini-0664:phyxminiproblems-problemphyxmini0664-ballastedtanksetup, def:physics:phyx-mini-0664:phyxminiproblems-problemphyxmini0664-matchesproblemandfiguredata, def:physics:phyx-mini-0664:phyxminiproblems-problemphyxmini0664-satisfiesprismaticdisplacementlaw}
  The 3 m² area and 10 m draft displace 30 m³ of ocean
\end{lemma}
\begin{proof}
  The conclusion follows from the governing relations and figure readouts named in the statement of displaced ocean volume at requested draft.
\end{proof}
\begin{lemma}[total supported mass at requested draft]
  \label{lem:physics:phyx-mini-0664:phyxminiproblems-problemphyxmini0664-total-supported-mass-at-requested-draft}
  \lean{PhyXMiniProblems.ProblemPhyXMini0664.total_supported_mass_at_requested_draft}
  \uses{def:physics:phyx-mini-0664:phyxminiproblems-problemphyxmini0664-ballastedtanksetup, def:physics:phyx-mini-0664:phyxminiproblems-problemphyxmini0664-totalsupportedmassinkilograms, def:physics:phyx-mini-0664:phyxminiproblems-problemphyxmini0664-matchesproblemandfiguredata, def:physics:phyx-mini-0664:phyxminiproblems-problemphyxmini0664-usesrecordedoceandensitycalibration, def:physics:phyx-mini-0664:phyxminiproblems-problemphyxmini0664-hasphysicalgravity, def:physics:phyx-mini-0664:phyxminiproblems-problemphyxmini0664-satisfiesprismaticdisplacementlaw, def:physics:phyx-mini-0664:phyxminiproblems-problemphyxmini0664-satisfiestotalweightlaw, def:physics:phyx-mini-0664:phyxminiproblems-problemphyxmini0664-satisfiesarchimedesprinciple, def:physics:phyx-mini-0664:phyxminiproblems-problemphyxmini0664-floatsinverticalstaticequilibrium}
  At the requested draft, buoyancy supports 997 × 30 = 29910 kg
\end{lemma}
\begin{proof}
  The conclusion follows from the governing relations and figure readouts named in the statement of total supported mass at requested draft.
\end{proof}
% --- Archon declaration topology end ---
% --- Archon physics formalization source end ---
